% Unit 110 English translation of chapter8.tex lines 760--924.
% Source: Methods of Algebra, Volume 2, commit
% 9a5803ff77dd3257484cb177f851a73770a59dd3 (CC BY 4.0).
% stable-unit-id: o014.aljabr2.chapter8.dold-kan-a
% Coverage: Section 8.5 from its opening through the Dold--Kan lemma for abelian groups;
% continued by Unit 111 at line 925.
% Editorial correction: the unused condition 0 <= k <= n in the source's
% definition of the normalized complex is omitted.

% segment-id: o014.li2.chapter8.dold-kan-a.h001
\section{The Dold--Kan Correspondence}\label{sec:Dold-Kan}
\hypertarget{o014.aljabr2.chapter8.dold-kan-a}{ }
% segment-id: o014.li2.chapter8.dold-kan-a.p001
Following the conventions of algebraic topology, for an additive category
$\mathcal{A}$ we consider the category of chain complexes over it
(Definition~\ref{def:complex}), denoted by $\cate{Ch}(\mathcal{A})$, together
with its subcategory
% segment-id: o014.li2.chapter8.dold-kan-a.q001
\[ \cate{Ch}_{\geq 0}(\mathcal{A}) := \left\{ A = (A_n, \partial_n)_{n \in \Z} : \text{chain complex}, \; \forall n < 0, \; A_n = 0 \right\}, \]
where $\partial_n=\partial_n^A:A_n\to A_{n-1}$. For an object of
$\cate{Ch}_{\geq0}(\mathcal{A})$, it plainly suffices to specify the data
$(A_n)_{n\geq0}$ and $(\partial_n)_{n\geq1}$. The category
$\cate{Ch}_{\geq m}(\mathcal{A})$ is defined similarly; compare the cochain
complex version in Definition~\ref{def:cplx-cat-variant}.
\index[sym1]{ChA@$\cate{Ch}(\mathcal{A})$}

% segment-id: o014.li2.chapter8.dold-kan-a.p002
By enlarging the chosen Grothendieck universe if necessary, we may assume that
$\mathcal{A}$ is a small additive category; see
\cite[Assumption 1.5.2]{Li1}. From $\mathcal{A}$ one can define the category
$\cate{s}\mathcal{A}$ of simplicial objects; the constant simplicial object
corresponding to $0$ is still denoted by
$0\in\Obj(\cate{s}\mathcal{A})$. This section aims to give a first account of
the relationships among these categories. More precisely, we shall define
three functors
% segment-id: o014.li2.chapter8.dold-kan-a.q002
\[\begin{tikzcd}
	\cate{Ch}_{\geq 0}(\mathcal{A}) \arrow[r, "\Gamma" description] & \cate{s}\mathcal{A} \arrow[bend left, l, "{\mathrm{N}}"] \arrow[bend right, l, "{\mathrm{C}}"']
\end{tikzcd}\]
where $\mathrm{N}$ is defined only when $\mathcal{A}$ is an abelian category.
Both $\mathrm{C}$ and $\mathrm{N}$ extend to the semisimplicial objects of
Remark~\ref{rem:semisimplicial}; in other words, neither of them involves the
degeneracy morphisms.

% segment-id: o014.li2.chapter8.dold-kan-a.p003
Our treatment follows \cite[\S 1.2]{LuHA}. We begin by making the relevant
definitions explicit.

% segment-id: o014.li2.chapter8.dold-kan-a.d001
\begin{definition}[Unnormalized chain complex]\label{def:unnormalized-chain}
	\index{unnormalized chain complex}
	\index{Moore chain complex}
	\index[sym1]{CX@$\mathrm{C}X$}
	For a semisimplicial object $X$ in $\mathcal{A}$ and every
	$n\in\Z_{\geq1}$, define
	% segment-id: o014.li2.chapter8.dold-kan-a.q003
	\[ \partial_n := \sum_{i=0}^n (-1)^i d_i: X_n \to X_{n-1}. \]
	The objects $(X_n)_{n\geq0}$ together with $(\partial_n)_n$ form an
	object $\mathrm{C}X$ of $\cate{Ch}_{\geq0}(\mathcal{A})$, called the
	\emph{unnormalized chain complex}, or \emph{Moore chain complex}, given
	by $X$. Thus $(\mathrm{C}X)_n=X_n$.
\end{definition}

% segment-id: o014.li2.chapter8.dold-kan-a.p004
We must show that $\partial_{n-1}\partial_n=0$ when $n>1$. Indeed, the
identity $i<j\implies d_id_j=d_{j-1}d_i$ in
\eqref{eqn:simplicial-identity} gives
% segment-id: o014.li2.chapter8.dold-kan-a.q004
\begin{multline*}
	\partial_{n-1} \partial_n = \sum_{i=0}^{n-1} \sum_{j=0}^n (-1)^{i+j} d_i d_j \\
	= \sum_{0 \leq i < j \leq n} (-1)^{i+j} d_{j-1} d_i + \sum_{n-1 \geq i \geq j \geq 0} (-1)^{i+j} d_i d_j = 0.
\end{multline*}
This construction is plainly canonical in $X$ and gives the functor
$\mathrm{C}$.

% segment-id: o014.li2.chapter8.dold-kan-a.eg001
\begin{example}\label{eg:classifying-space-std-cplx}
	Let $\Bbbk$ be a commutative ring and
	$\mathcal{A}=\Bbbk\dcate{Mod}$. For a group $\Gamma$,
	Example~\ref{eg:classifying-space} gives the corresponding simplicial set
	$\mathrm{E}\Gamma$. Define a simplicial object
	$\Bbbk\mathrm{E}\Gamma$ in $\Bbbk\dcate{Mod}$ whose term in degree $n$
	is the free $\Bbbk$-module with basis $(\mathrm{E}\Gamma)_n$, and whose
	$d_i$ and $s_i$ are extended linearly. We claim that
	$\mathrm{C}(\Bbbk\mathrm{E}\Gamma)$ is precisely the chain complex
	$\mathsf{L}:=(\mathsf{L}_n,\partial'_n)_n$ introduced in
	Definition~\ref{def:resolution-trivial-module}. First,
	$\mathsf{L}_n$ is by definition the free $\Bbbk$-module with basis
	$\Gamma^{n+1}$. Take the map from
	$(\mathrm{E}\Gamma)_n=\Gamma^{n+1}$ to $\mathsf{L}_n$
	% segment-id: o014.li2.chapter8.dold-kan-a.q005
	\[ \left( g_1, \ldots, g_{n+1} \right) \mapsto \left(g_1 \cdots g_{n+1}, \; \ldots , \; g_n g_{n+1}, \; g_{n+1}\right). \]
	Extend it linearly to $(\Bbbk\mathrm{E}\Gamma)_n$. The reader may verify
	that $\sum_i(-1)^id_i$ thereby corresponds to
	$\partial'_n:\mathsf{L}_n\to\mathsf{L}_{n-1}$, giving the asserted
	isomorphism.

	Moreover, $\Gamma$ acts on the right on each $(\mathrm{E}\Gamma)_n$ by
	$(g_1,\ldots,g_{n+1})\xmapsto{g}(g_1,\ldots,g_n,g_{n+1}g)$. In the basis
	coordinates of $\mathsf{L}_n$, this action becomes
	$(g_0,\ldots,g_n)\xmapsto{g}(g_0g,\ldots,g_ng)$. Thus we recover the
	right $\Gamma$-action on $\mathsf{L}$ from
	Definition~\ref{def:resolution-trivial-module}. Consequently, the
	complexes obtained by taking degreewise $\Gamma$-coinvariants of the
	chain complexes $\mathrm{C}(\Bbbk\mathrm{E}\Gamma)$ and $\mathsf{L}$
	are also isomorphic. On the $\mathrm{C}(\Bbbk\mathrm{E}\Gamma)$ side,
	this amounts to taking the degreewise quotient of $\mathrm{E}\Gamma$ and
	then the free $\Bbbk$-module. It follows that
	% segment-id: o014.li2.chapter8.dold-kan-a.q006
	\[ \mathrm{C}(\Bbbk\mathrm{B}\Gamma) \simeq \mathsf{L} \dotimes{\Bbbk[\Gamma]} \Bbbk, \]
	where the tensor product involves the augmentation homomorphism
	$\Bbbk[\Gamma]\to\Bbbk$. As a consequence, group homology
	(Definition~\ref{def:group-coh-ho}) is interpreted as
	% segment-id: o014.li2.chapter8.dold-kan-a.q007
	\[ \Hm_n(\mathrm{B}\Gamma; \Bbbk) := \Hm_n \mathrm{C}(\Bbbk\mathrm{B}\Gamma) \simeq \Hm_n(\Gamma, \Bbbk), \quad n \in \Z_{\geq 0}. \]

	The left-hand side features the homology of a simplicial set, a general
	construction to be introduced in
	Definition~\ref{def:simplicial-set-homology}. The exercises in this
	chapter will give a similar isomorphism for the normalized version
	$\overline{\mathsf{L}}$ of $\mathsf{L}$, corresponding to the normalized
	chain complex $\mathrm{N}(\Bbbk\mathrm{E}\Gamma)$ to be discussed in
	Definition~\ref{def:normalized-cplx}.
\end{example}

% segment-id: o014.li2.chapter8.dold-kan-a.d002
\begin{definition}\label{def:DK-Gamma}
	Let $(A_n,\partial_n)_{n\geq0}$ be an object of
	$\cate{Ch}_{\geq0}(\mathcal{A})$. Define an object $\Gamma(A)$ of
	$\cate{s}\mathcal{A}$ as follows:
	% segment-id: o014.li2.chapter8.dold-kan-a.q008
	\[ (\Gamma A)_n := \bigoplus_{t: [n] \twoheadrightarrow [k]} A_k. \]

	For every morphism $f:[m]\to[n]$ in $\simpDelta$, the corresponding
	morphism $f^*:(\Gamma A)_n\to(\Gamma A)_m$ is represented, with respect
	to the direct-sum decomposition above, by the matrix
	% segment-id: o014.li2.chapter8.dold-kan-a.q009
	\[ \left( \Phi_{u, t} : A_k \to A_l \right)_{\substack{ u: [m] \twoheadrightarrow [l] \\ t: [n] \twoheadrightarrow [k] }}. \]
	This matrix is determined as follows. Given
	$t:[n]\twoheadrightarrow[k]$, take the surjection--injection
	factorization of $tf$,
	% segment-id: o014.li2.chapter8.dold-kan-a.q010
	\begin{equation}\label{eqn:Dold-Kan-epi-mono}\begin{tikzcd}[row sep=tiny]
		& {[n]} \arrow[twoheadrightarrow, rd, "t"] & \\
		{[m]} \arrow[ru, "f"] \arrow[twoheadrightarrow, rd, "u"'] & & {[k]} \\
		& {[l]} \arrow[hookrightarrow, ru, "v"'] &
	\end{tikzcd}\end{equation}
	\begin{itemize}
		\item if $l=k$, set $\Phi_{u,t}=\identity$;
		\item if $l=k-1$ and $v=\mathrm{d}^0$ (the order-preserving
		injection omitting $0$), set
		$\Phi_{u,t}=\partial_k:A_k\to A_{k-1}$;
		\item in all other cases, set $\Phi_{u,t}=0$.
	\end{itemize}
	Notice that once $f$ and $t$ are given, there is at most one $u$ for which
	$\Phi_{u,t}\neq0$.
\end{definition}

% segment-id: o014.li2.chapter8.dold-kan-a.p005
The verification that $\Gamma A\in\Obj(\cate{s}\mathcal{A})$ is somewhat
tedious and is omitted. The heart of the construction lies in the direct
summand $A_n$ of $(\Gamma A)_n$ and its behavior under order-preserving
injections. All lower-degree direct summands $A_k$ arise from degeneracies
induced by surjections $[n]\twoheadrightarrow[k]$; the definition of $f^*$
in the general case merely implements this idea.

% segment-id: o014.li2.chapter8.dold-kan-a.p006
The construction above is canonical and gives a functor
$\Gamma:\cate{Ch}_{\geq0}(\mathcal{A})\to\cate{s}\mathcal{A}$.

% segment-id: o014.li2.chapter8.dold-kan-a.d003
\begin{definition}\label{def:normalized-cplx}
	\index{normalized chain complex}
	\index[sym1]{NX@$\mathrm{N}X$}
	Let $\mathcal{A}$ be an abelian category. For a semisimplicial object $X$
	in $\mathcal{A}$ and every $n\in\Z_{\geq0}$, define
	% segment-id: o014.li2.chapter8.dold-kan-a.q011
	\begin{align*}
		(\mathrm{N} X)_n & := \bigcap_{i=1}^n \Ker\left[ d_i: X_n \to X_{n-1} \right], \quad n \geq 1, \\
		(\mathrm{N} X)_0 & := X_0.
	\end{align*}

	The morphisms $\mathrm{N}X_n\to\mathrm{N}X_{n-1}$ induced by
	$X_n\xrightarrow{d_0}X_{n-1}$ (verify this using
	\eqref{eqn:simplicial-identity}) are henceforth denoted by $\partial_n$.
	Together with these objects, they form an object $\mathrm{N}X$ of
	$\cate{Ch}_{\geq0}(\mathcal{A})$, called the \emph{normalized chain
	complex} corresponding to $X$.
\end{definition}

% segment-id: o014.li2.chapter8.dold-kan-a.pr001
\begin{proposition}\label{prop:Dold-Kan-u}
	For $X$ as above and every $n\geq0$, write
	$u_n:(\mathrm{N}X)_n\hookrightarrow X_n$ for the natural inclusion. Then
	$(u_n)_{n\geq0}$ forms a monomorphism
	$u:\mathrm{N}X\to\mathrm{C}X$ in
	$\cate{Ch}_{\geq0}(\mathcal{A})$; this construction is functorial in $X$.
\end{proposition}
% segment-id: o014.li2.chapter8.dold-kan-a.pf001
\begin{proof}
	This follows directly from the definitions of $\mathrm{N}X$ and
	$\mathrm{C}X$.
\end{proof}

% segment-id: o014.li2.chapter8.dold-kan-a.p007
The next few results are preparations for the Dold--Kan theorem,
Theorem~\ref{prop:Dold-Kan}.

% segment-id: o014.li2.chapter8.dold-kan-a.le001
\begin{lemma}\label{prop:DK-N-unit}
	Let $\mathcal{A}$ be an abelian category. For every
	$A\in\Obj(\cate{Ch}_{\geq0}(\mathcal{A}))$ and $n\geq0$, we have
	$\mathrm{N}\Gamma(A)_n=A_n$. These identifications give an isomorphism
	of functors
	$\eta:\identity_{\cate{Ch}_{\geq0}(\mathcal{A})}\rightiso
	\mathrm{N}\Gamma$.
\end{lemma}
% segment-id: o014.li2.chapter8.dold-kan-a.pf002
\begin{proof}
	Given $t:[n]\twoheadrightarrow[k]$ with $k<n$, one can always choose
	$1\leq i\leq n$ such that $t^{-1}(t(i))$ has at least two elements.
	Consequently $t\mathrm{d}^i$ is surjective, and the following diagram
	commutes:
	% segment-id: o014.li2.chapter8.dold-kan-a.q012
	\[\begin{tikzcd}[row sep=tiny]
		& {[n]} \arrow[twoheadrightarrow, rd, "t"] & \\
		{[n-1]} \arrow[ru, "{\mathrm{d}^i}"] \arrow[twoheadrightarrow, rd, "{t\mathrm{d}^i}"'] & & {[k]} \\
		& {[k]} \arrow[ru, "{\identity}"'] &
	\end{tikzcd}\]
	Definition~\ref{def:DK-Gamma} then shows that
	$\mathrm{N}\Gamma(A)_n$ is contained in the direct summand $A_n$
	corresponding to $\identity:[n]\to[n]$. The same definition readily
	gives $\mathrm{N}\Gamma(A)_n=A_n$ and shows that
	$\mathrm{N}\Gamma(A)_{n+1}\to\mathrm{N}\Gamma(A)_n$ is precisely
	$\partial_{n+1}:A_{n+1}\to A_n$. This proves the result.
\end{proof}

% segment-id: o014.li2.chapter8.dold-kan-a.le002
\begin{lemma}\label{prop:GammaN-adjunction}
	Let $\mathcal{A}$ be an abelian category. The functors above form an
	adjunction
	% segment-id: o014.li2.chapter8.dold-kan-a.q013
	\[\begin{tikzcd}
		\Gamma: \cate{Ch}_{\geq 0}(\mathcal{A}) \arrow[shift left, r] & \cate{s}\mathcal{A} : \mathrm{N} \arrow[shift left, l].
	\end{tikzcd}\]
	\begin{itemize}
		\item Its unit morphism is the isomorphism $\eta$ of
		Lemma~\ref{prop:DK-N-unit}.
		\item Its counit morphism
		$\epsilon:\Gamma\mathrm{N}\to\identity_{\cate{s}\mathcal{A}}$ is
		described as follows. For $t:[n]\twoheadrightarrow[k]$, on the
		corresponding direct summand $(\mathrm{N}X)_k$ of
		$\Gamma\mathrm{N}(X)_n$, the map $\epsilon_{X,n}$ is the composite
		$(\mathrm{N}X)_k\subset X_k\xrightarrow{t^*}X_n$.
	\end{itemize}
\end{lemma}
% segment-id: o014.li2.chapter8.dold-kan-a.pf003
\begin{proof}
	For $A\in\Obj(\cate{Ch}_{\geq0}(\mathcal{A}))$ and
	$X\in\Obj(\cate{s}\mathcal{A})$, we first prove that the composite
	% segment-id: o014.li2.chapter8.dold-kan-a.q014
	\[ \Hom_{\cate{s}\mathcal{A}}(\Gamma(A), X) \xrightarrow{\mathrm{N}} \Hom_{\cate{Ch}_{\geq 0}(\mathcal{A})}(\mathrm{N}\Gamma(A), \mathrm{N}(X)) \xrightarrow{\eta_A^*} \Hom_{\cate{Ch}_{\geq 0}(\mathcal{A})}(A, \mathrm{N}(X)) \]
	is a bijection. Its inverse is defined concretely as follows. Given
	$\phi=(\phi_n)_{n\geq0}:A\to\mathrm{N}(X)$, define a morphism
	% segment-id: o014.li2.chapter8.dold-kan-a.q015
	\[ \Phi_n: (\Gamma A)_n = \bigoplus_{t: [n] \twoheadrightarrow [k]} A_k \to X_n \]
	whose restriction to the direct summand corresponding to
	$t:[n]\twoheadrightarrow[k]$ is the composite
	% segment-id: o014.li2.chapter8.dold-kan-a.q016
	\[ A_k \xrightarrow{\phi_k} (\mathrm{N}X)_k \subset X_k \xrightarrow{t^*} X_n. \]
	For such a $t$ and an arbitrary $f:[m]\to[n]$, take the
	surjection--injection factorization of $tf$ as in
	\eqref{eqn:Dold-Kan-epi-mono}, and consider the following diagram in
	$\mathcal{A}$ (see Definition~\ref{def:DK-Gamma}):
	% segment-id: o014.li2.chapter8.dold-kan-a.q017
	\[\begin{tikzcd}[column sep=large]
		A_k \arrow[r, "\phi_k"] \arrow[d, "{\Phi_{u, t}}"'] & (\mathrm{N}X)_k \arrow[hookrightarrow, r] \arrow[d, "{v^*|_{(\mathrm{N}X)_k}}"] & X_k \arrow[r, "{t^*}"] \arrow[d, "{v^*}"] & X_n \arrow[d, "{f^*}"] \\
		A_l \arrow[r, "\phi_l"'] & (\mathrm{N}X)_l \arrow[hookrightarrow, r] & X_l \arrow[r, "{u^*}"'] & X_m.
	\end{tikzcd}\]
	The left square commutes by the definition of the normalized chain complex
	$\mathrm{N}X$, the middle square plainly commutes, and the right square
	commutes because $tf=vu$. Thus $\Phi=(\Phi_n)_n$ is a morphism in
	$\cate{s}\mathcal{A}$.

	The verification that $\eta_A^*\mathrm{N}$ and $\phi\mapsto\Phi$ are
	mutually inverse is routine. In the description of $\phi\mapsto\Phi$,
	take $\phi=\identity_{\mathrm{N}X}$; the result is exactly the asserted
	$\epsilon$.
\end{proof}

% segment-id: o014.li2.chapter8.dold-kan-a.le003
\begin{lemma}\label{prop:DK-Ab}
	Take $\mathcal{A}=\cate{Ab}$. Then $\Gamma$ is an equivalence; more
	precisely, the adjunction
	% segment-id: o014.li2.chapter8.dold-kan-a.q018
	$\begin{tikzcd}
		\Gamma: \cate{Ch}_{\geq 0}(\cate{Ab}) \arrow[shift left, r] & \cate{s}\cate{Ab} : \mathrm{N} \arrow[shift left, l]
	\end{tikzcd}$
	is an adjoint equivalence in the sense of
	\cite[Theorem 2.6.12]{Li1}.
\end{lemma}
% segment-id: o014.li2.chapter8.dold-kan-a.pf004
\begin{proof}
	It suffices to prove that the counit morphism $\epsilon$ of
	Lemma~\ref{prop:GammaN-adjunction} is an isomorphism. Let
	$X\in\Obj(\cate{s}\mathcal{A})$. We claim that
	$\epsilon_{X,n}:\Gamma(\mathrm{N}X)_n\to X_n$ is injective for every
	$n\geq0$. Let $x=(x_t)_t$ be an element of its domain, where $t$ ranges
	over all order-preserving surjections $[n]\twoheadrightarrow[k]$. For a
	given $t$, define the order-preserving injection
	% segment-id: o014.li2.chapter8.dold-kan-a.q019
	\begin{align*}
		s = s_t : [k] & \hookrightarrow [n] \\
		i & \mapsto \min t^{-1}(i),
	\end{align*}
	which satisfies $ts=\identity_{[k]}$. Suppose that $x\neq0$, and write
	% segment-id: o014.li2.chapter8.dold-kan-a.q020
	\[ S := \left\{t: x_t \neq 0 \right\} \neq \emptyset. \]
	Choose the smallest $k$ for which there is a
	$t:[n]\twoheadrightarrow[k]$ in $S$, then choose such a $t$ minimizing
	$\sum_{i=0}^k\min t^{-1}(i)$, and construct $s$. We shall prove that
	$s^*(\epsilon_{X,n}(x))=x_t\in X_k$, which shows that
	$\epsilon_{X,n}(x)\neq0$.

	By the concrete description of $\epsilon_{X,n}$, it is enough to show
	that for every $t':[n]\twoheadrightarrow[k']$,
	$s^*(t')^*(x_{t'})\neq0$ implies $t=t'$. Set
	$u:=t's:[k]\to[k']$. If $t'\notin S$, then $x_{t'}=0$. Thus it remains
	to consider $t'\in S$ satisfying
	$s^*(t')^*(x_{t'})=u^*(x_{t'})\neq0$.

	Since $t$ is order-preserving and surjective, $\min t^{-1}(0)=0$, and
	the same holds for $t'$, so $u(0)=0$. Moreover,
	$x_{t'}\in(\mathrm{N}X)_{k'}$ implies that
	$u^*(x_{t'})\neq0$ is possible only if
	$\Image(u)\supset\{1,\ldots,k'\}$. Hence we may assume that $u$ is
	surjective. The minimality of $k$ then forces $k'=k$ and
	$u=\identity$.

	Consequently, $t'(\min t^{-1}(i))=i$ for $i=0,\ldots,k$, and hence
	$\min(t')^{-1}(i)\leq\min t^{-1}(i)$. By the choice of $t$, we obtain
	$\min(t')^{-1}(i)=\min t^{-1}(i)$ for every $0\leq i\leq k$.
	An order-preserving surjection is determined by the initial points of its
	fibers, so this is equivalent to $t=t'$. Thus $\epsilon_{X,n}$ is
	injective.

	We now prove by induction on $n$ that $\epsilon_{X,n}$ is surjective. We
	claim that, for every $0\leq i\leq n$,
	% segment-id: o014.li2.chapter8.dold-kan-a.q021
	\[ \Image\left( \epsilon_{X, n} \right) \supset X(i)_n := \bigcap_{i < j \leq n} \Ker(d_j) \subset X_n. \]

	For $i=0$, we have $X(i)_n=(\mathrm{N}X)_n$, and the inclusion above
	follows readily from the concrete description of $\epsilon_{X,n}$. Our
	goal is the case $i=n$. Let $n\geq i\geq1$ and $y\in X(i)_n$. The
	induction hypothesis for $n-1$ and the description of $\epsilon_{X,n}$
	give
	% segment-id: o014.li2.chapter8.dold-kan-a.q022
	\[ s_{i-1} d_i(y) \in s_{i-1}\left( \Image\left(\epsilon_{X, n-1}\right) \right) \subset \Image\left( \epsilon_{X, n} \right). \]

	On the other hand, the simplicial identities
	\eqref{eqn:simplicial-identity} give
	% segment-id: o014.li2.chapter8.dold-kan-a.q023
	\begin{gather*}
		d_i s_{i-1} d_i = d_i, \\
		j > i \implies d_j s_{i-1} d_i = s_{i-1} d_{j-1} d_i = s_{i-1} d_i d_j.
	\end{gather*}
	Thus $y-s_{i-1}d_i(y)\in X(i-1)_n$, and the induction hypothesis for
	$i-1$ shows that this element lies in
	$\Image(\epsilon_{X,n})$. Therefore
	$y\in\Image(\epsilon_{X,n})$, as required.
\end{proof}
